% GENERATED FILE. Edit canonical lessons, then run npm run generate:books.
% canonical-source-hash: fnv1a64:e94dc42a77b79b1b
% canonical-lessons: FA-C489-moaddab, FA-C489-biadab, FA-C489-zudranj, FA-C489-kamru, FA-C489-maskhare-kardan

\chapter{Manners and Temperament}
\label{ch:fa-manners-and-temperament}
\hlchaptermodality{489}

\begin{chapteropening}
\textbf{By the end of this chapter:} \emph{I can describe manners and temperament.}

Complete the last lesson of chapter 489: I can describe manners and temperament.
\end{chapteropening}

\section[\texorpdfstring{\fa{مودب} (mo'addab)}{مودب (mo'addab)}]{\fa{مودب} (mo'addab) — polite, well-mannered}

\label{lesson:FA-C489-moaddab}



\begin{quote}
Before the new one: say the Persian for generous, then the Persian for caring.
\end{quote}

\subsection*{You'll want to know: \fa{مودب}}
\textbf{\fa{مودب}} (\emph{mo'addab}) — \textquotedblleft{}polite, well-mannered\textquotedblright{}.

\begin{quote}
\fa{بچه} \fa{های} \fa{آنها} \fa{خیلی} \fa{مودب} \fa{هستند}.
\end{quote}

\begin{quote}
\emph{bachchehâ-ye ânhâ kheyli mo'addab hastand.}
\end{quote}

\begin{quote}
Their children are very polite.
\end{quote}

Built on \emph{adab} 'good manners'. The opposite is \emph{biadab}.

\begin{cousinweb}
\emph{mo'addab} is a loan, usually traced to Arabic \emph{muʾaddab}, 'well-bred', a word built on \emph{adab}, 'manners'. EtymDB has no entry for it.
\end{cousinweb}

\begin{grammarlens}[title={What you've built}]
One more word for telling, explaining and describing at B1.
\end{grammarlens}

\subsection*{Your turn}
\begin{itemize}
\raggedright
  \item \emph{Say it:} \emph{mo'addab}
  \item \emph{Say it:} \emph{mo'addab}, once more
  \item \emph{Say it:} a sentence of your own with \emph{mo'addab}
  \item \emph{Recall:} say the Persian for generous, then the Persian for caring, then say \emph{mo'addab} again
\end{itemize}

\begin{tcolorbox}[breakable,colback=teal!4,colframe=teal!35!black,title={Before you move on}]
What does \fa{مودب} mean? (\textquotedblleft{}Polite, well-mannered\textquotedblright{}.) Say it once more.
\end{tcolorbox}
\section[\texorpdfstring{\fa{بیادب} (biadab)}{بیادب (biadab)}]{\fa{بیادب} (biadab) — rude, impolite}

\label{lesson:FA-C489-biadab}



\begin{quote}
Before the new one: say the Persian for caring, then the Persian for polite.
\end{quote}

\subsection*{You'll want to know: \fa{بیادب}}
\textbf{\fa{بیادب}} (\emph{biadab}) — \textquotedblleft{}rude, impolite\textquotedblright{}.

\begin{quote}
\fa{آن} \fa{مرد} \fa{با} \fa{فروشنده} \fa{بیادب} \fa{بود}.
\end{quote}

\begin{quote}
\emph{ân mard bâ forushande biadab bud.}
\end{quote}

\begin{quote}
That man was rude to the shop assistant.
\end{quote}

\emph{bi-} 'without' + \emph{adab} 'manners'. \emph{bâ kasi biadab budan} — to be rude to someone.

\begin{cousinweb}
\emph{biadab} is built in Persian from the prefix \emph{bi-}, 'without', and \emph{adab}, 'manners', a loan usually traced to Arabic \emph{adab}. EtymDB has no entry for it.
\end{cousinweb}

\begin{grammarlens}[title={What you've built}]
One more word for telling, explaining and describing at B1.
\end{grammarlens}

\subsection*{Your turn}
\begin{itemize}
\raggedright
  \item \emph{Say it:} \emph{biadab}
  \item \emph{Say it:} \emph{biadab}, once more
  \item \emph{Say it:} a sentence of your own with \emph{biadab}
  \item \emph{Recall:} say the Persian for caring, then the Persian for polite, then say \emph{biadab} again
\end{itemize}

\begin{tcolorbox}[breakable,colback=teal!4,colframe=teal!35!black,title={Before you move on}]
What does \fa{بیادب} mean? (\textquotedblleft{}Rude, impolite\textquotedblright{}.) Say it once more.
\end{tcolorbox}
\section[\texorpdfstring{\fa{زودرنج} (zudranj)}{زودرنج (zudranj)}]{\fa{زودرنج} (zudranj) — touchy, easily hurt}

\label{lesson:FA-C489-zudranj}



\begin{quote}
Before the new one: say the Persian for polite, then the Persian for rude.
\end{quote}

\subsection*{You'll want to know: \fa{زودرنج}}
\textbf{\fa{زودرنج}} (\emph{zudranj}) — \textquotedblleft{}touchy, easily hurt\textquotedblright{}.

\begin{quote}
\fa{او} \fa{زودرنج} \fa{است} \fa{و} \fa{با} \fa{کوچکترین} \fa{حرفی} \fa{ناراحت} \fa{میشود}.
\end{quote}

\begin{quote}
\emph{u zudranj ast va bâ kuchektarin harfi nârâhat mishavad.}
\end{quote}

\begin{quote}
She's touchy and gets upset at the smallest remark.
\end{quote}

\emph{zud} 'quickly' + \emph{ranj} 'hurt', the noun behind \emph{ranjidan}.

\begin{cousinweb}
\emph{zudranj} is built on two native Persian words: \emph{zud}, 'quickly', and \emph{ranj}, 'suffering', the noun behind \emph{ranjidan}: quick to be hurt. EtymDB has no entry for it.
\end{cousinweb}

\begin{grammarlens}[title={What you've built}]
One more word for telling, explaining and describing at B1.
\end{grammarlens}

\subsection*{Your turn}
\begin{itemize}
\raggedright
  \item \emph{Say it:} \emph{zudranj}
  \item \emph{Say it:} \emph{zudranj}, once more
  \item \emph{Say it:} a sentence of your own with \emph{zudranj}
  \item \emph{Recall:} say the Persian for polite, then the Persian for rude, then say \emph{zudranj} again
\end{itemize}

\begin{tcolorbox}[breakable,colback=teal!4,colframe=teal!35!black,title={Before you move on}]
What does \fa{زودرنج} mean? (\textquotedblleft{}Touchy, easily hurt\textquotedblright{}.) Say it once more.
\end{tcolorbox}
\section[\texorpdfstring{\fa{کمرو} (kamru)}{کمرو (kamru)}]{\fa{کمرو} (kamru) — shy, timid (in company)}

\label{lesson:FA-C489-kamru}



\begin{quote}
Before the new one: say the Persian for rude, then the Persian for touchy.
\end{quote}

\subsection*{You'll want to know: \fa{کمرو}}
\textbf{\fa{کمرو}} (\emph{kamru}) — \textquotedblleft{}shy, timid (in company)\textquotedblright{}.

\begin{quote}
\fa{پسرشان} \fa{کمرو} \fa{است} \fa{و} \fa{در} \fa{مهمانی} \fa{حرف} \fa{نمیزند}.
\end{quote}

\begin{quote}
\emph{pesarashân kamru ast va dar mehmâni harf nemizanad.}
\end{quote}

\begin{quote}
Their son is shy and doesn't talk at parties.
\end{quote}

\emph{kam} 'little' + \emph{ru} 'face': little face, little boldness. The opposite is \emph{porru} 'cheeky'.

\begin{cousinweb}
\emph{kamru} is built on two native Persian words, \emph{kam}, 'little', and \emph{ru}, 'face'. It is the mirror of \emph{porru}, 'cheeky', built on \emph{por}, 'full'.
\end{cousinweb}

\begin{grammarlens}[title={What you've built}]
One more word for telling, explaining and describing at B1.
\end{grammarlens}

\subsection*{Your turn}
\begin{itemize}
\raggedright
  \item \emph{Say it:} \emph{kamru}
  \item \emph{Say it:} \emph{kamru}, once more
  \item \emph{Say it:} a sentence of your own with \emph{kamru}
  \item \emph{Recall:} say the Persian for rude, then the Persian for touchy, then say \emph{kamru} again
\end{itemize}

\begin{tcolorbox}[breakable,colback=teal!4,colframe=teal!35!black,title={Before you move on}]
What does \fa{کمرو} mean? (\textquotedblleft{}Shy, timid (in company)\textquotedblright{}.) Say it once more.
\end{tcolorbox}
\section[\texorpdfstring{\fa{مسخره} \fa{کردن} (maskhare kardan)}{مسخره کردن (maskhare kardan)}]{\fa{مسخره} \fa{کردن} (maskhare kardan) — to make fun of, to mock}

\label{lesson:FA-C489-maskhare-kardan}



\begin{quote}
Before the new one: say the Persian for touchy, then the Persian for shy.
\end{quote}

\subsection*{You'll want to know: \fa{مسخره} \fa{کردن}}
\textbf{\fa{مسخره} \fa{کردن}} (\emph{maskhare kardan}) — \textquotedblleft{}to make fun of, to mock\textquotedblright{}.

\begin{quote}
\fa{او} \fa{لباس} \fa{من} \fa{را} \fa{مسخره} \fa{کرد}.
\end{quote}

\begin{quote}
\emph{u lebâs-e man râ maskhare kard.}
\end{quote}

\begin{quote}
He made fun of my clothes.
\end{quote}

\emph{maskhare} 'ridiculous; a laughing stock' + \emph{kardan}; present stem \emph{kon-}. The person mocked takes \emph{râ}.

\begin{cousinweb}
\emph{maskhare} comes from Arabic \emph{masḫara}, 'buffoon, anything ridiculous' (EtymDB). \emph{kardan}, 'to do', is inherited Persian.
\end{cousinweb}

\begin{grammarlens}[title={What you've built}]
One more word for telling, explaining and describing at B1.
\end{grammarlens}

\subsection*{Your turn}
\begin{itemize}
\raggedright
  \item \emph{Say it:} \emph{maskhare kardan}
  \item \emph{Say it:} \emph{maskhare kardan}, once more
  \item \emph{Say it:} a sentence of your own with \emph{maskhare kardan}
  \item \emph{Recall:} say the Persian for touchy, then the Persian for shy, then say \emph{maskhare kardan} again
\end{itemize}

\begin{tcolorbox}[breakable,colback=teal!4,colframe=teal!35!black,title={Before you move on}]
What does \fa{مسخره} \fa{کردن} mean? (\textquotedblleft{}To make fun of, to mock\textquotedblright{}.) Say it once more.
\end{tcolorbox}
